% Appendix, sections A through N.

% ============================================================================
% A — The Prior Routing Audit
% ============================================================================
\section{The Prior Routing Audit}
\label{apx:audit}

The cost-routing premise of \S\ref{sec:problem} rests on a zero-cost replay of
published per-task outcome matrices.  Table~\ref{tab:apx-audit} records the
structural statistics behind that replay: set containment, the unique-solver
shell, and the gap between the best learned router and always picking the
strongest model, across three benchmarks and matched four-model pools.

\begin{table}[H]
\centering
\footnotesize
\caption{\textbf{Why we do not route for accuracy.}
Zero-cost audit of published per-task outcome matrices on matched four-model
pools.  \emph{Containment} measures solve-set nesting; the
\emph{unique-solver shell} is the share of routable tasks only one model
solves.  On all three benchmarks, no learned router beat always calling the
strongest model.  This is background motivation, not a system result.}
\label{tab:apx-audit}
% ===========================================================================
% tab-apx-audit.tex — APPENDIX. The zero-cost audit of public per-task
% outcome matrices that motivated this work.
% BODY ONLY (tabular). Wrapper/caption/label live in the appendix file.
%
% PROVENANCE — every cell.
%   Source: internal records, "Master table (structural stats
%   at matched 4-model realistic pools; routing deltas pooled over LORO)".
%   Column order in the source is Verified | Multilingual | Pro. Verified
%   verbatim:
%     mean containment (Guttman)          0.941 | 0.912 | 0.773
%     unique-solver shell (% of routable) 21.5% | 20.8% | 30.8%
%     best router - always-strongest      -0.006 (p=.66) ns | +0.010 (p=.24)
%                                         ns | -0.004 (p=.90) ns
%   Supporting rows from the same table, used in the caption:
%     best single model 0.768 | 0.726 | 0.449
%   Source's own reading, verbatim: "the harder the benchmark, the LESS
%     nested ... and the bigger the genuinely-complementary shell"; and the
%     decisive line, "But — the decisive twist — routing still does not win on
%     Pro. Best router (ours-KNN) ties always-strongest (-0.004, p=0.90) ...
%     always-strongest is unbeaten on all three matrices."
%   Comparability note the source insists on: containment and shell are
%     computed on identical 4-model pools so the model-count confound is
%     removed ("mean-over-all-4-subsets containment = 0.934 / 0.888 / 0.773").
%
%   FRAMING RULE (BINDING): "accuracy-routing-beats-always-best = background only (v1
%   finding: impossible without specialist pools)". This table is background
%   motivation and nothing else. It must NOT be presented as a contribution,
%   a baseline comparison, or a result of this system: it is a prior,
%   zero-cost audit of published outcome matrices that explains why this work
%   routes for COST at matched accuracy rather than for accuracy.
% ===========================================================================
\begin{tabular}{@{}lrrr@{}}
\toprule
 & Verified & Multiling. & Pro \\
\midrule
Best single model      & 0.768 & 0.726 & 0.449 \\
Mean containment       & 0.941 & 0.912 & 0.773 \\
Unique-solver shell    & 21.5\% & 20.8\% & 30.8\% \\
\addlinespace
Best router $-$ best model & $-0.006$ & $+0.010$ & $-0.004$ \\
\quad $p$              & 0.66 & 0.24 & 0.90 \\
\bottomrule
\end{tabular}

\end{table}



% ============================================================================
% B — Router Threshold Sweep
% ============================================================================
\section{Router Threshold Sweep}
\label{apx:theta-sweep}

Table~\ref{tab:apx-theta-sweep} sweeps the routing threshold $\theta$ that
the calibration's matched-point rule selects.  At $\theta{=}0.25$ every task
clears the cheapest fixer's gate and the router degenerates into the no-router
ablation: 159 solves at \$0.227 per solve, identical to always calling
\fxkimi.  Raising the threshold to the operating point $\theta{=}0.30$ diverts
three tasks to \fxflash at a marginal cost of \$0.003 per solve, without
changing the solve count.  Above the operating point, the router's growing
caution pushes progressively more tasks onto expensive fixers: cost per task
rises monotonically while the solve count ceases to be exactly measurable,
because the diverted tasks lack a with-handoff outcome on their new fixer.
We report those rows as bracketing intervals rather than point estimates (the
table note details the convention).  On this benchmark, routing collapses to
a cost-allocation decision: no threshold examined buys additional solves
beyond the 159 that the cheapest fixer already delivers.

\begin{table}[H]
\centering
\footnotesize
\caption{\textbf{The gate threshold sweep.}
Raising $\theta$ pushes tasks off the cheapest fixer onto more expensive ones:
cost per task rises monotonically while the solve count ceases to be exactly
measurable, since diverted tasks lack a with-handoff outcome on their new
fixer.  The operating point $\theta{=}0.30$ routes all but three tasks to
\fxkimi and is the last threshold at which every number is exact.}
\label{tab:apx-theta-sweep}
% ===========================================================================
% tab-apx-theta-sweep.tex — APPENDIX. Gate threshold sweep: what the router
% does to routing, solves and cost as theta moves.
% BODY ONLY (tabular). Wrapper/caption/label live in the appendix file; the
% appendix wrapper supplies \footnotesize, so this body sets no font size,
% only a tighter \tabcolsep (eight columns in one 3.5in text column).
% Provenance: code/router/router_spec.json, data/receipts/outcomes_per_task.json.
%
% Convention: exact rows (0.25, 0.30) print a point solve count; bounded rows
%   print the bracketing interval [lo, hi] and NEVER a bare midpoint, here or
%   in caption or prose. The star marks a $/solve computed from a solo-filled
%   point estimate of the solve count; $/task carries no star because it is
%   fixed by the routing split. theta = 0.30 is the operating point. Above the
%   operating point, tasks passing no gate go to the most capable fixer
%   (stated in the note); at 0.25/0.30 the fall-through never fires.
% ===========================================================================
\begingroup\setlength{\tabcolsep}{2pt}%
\begin{tabular}{@{}lrrrrrrr@{}}
\toprule
& \multicolumn{4}{c}{Routing split (tasks)} & & & \\
\cmidrule(lr){2-5}
$\theta$
  & \parbox[b]{0.42in}{\raggedleft\fxkimi}
  & \parbox[b]{0.42in}{\raggedleft\fxflash}
  & \parbox[b]{0.42in}{\raggedleft\fxgpt}
  & \parbox[b]{0.42in}{\raggedleft\fxopus}
  & Solves & \$/task & \$/solve \\
\midrule
$0.25$              & 266 &  0 &  0 &  0 & 159        & \$0.136 & \$0.227 \\
$0.30^{\dagger}$    & 263 &  3 &  0 &  0 & 159        & \$0.137 & \$0.230 \\
$0.35$              & 236 & 19 & 10 &  1 & [149, 169] & \$0.154 & \$0.272$^{*}$ \\
$0.40$              & 203 & 34 & 28 &  1 & [132, 179] & \$0.178 & \$0.337$^{*}$ \\
$0.50$              & 124 & 34 & 72 & 36 & [94, 196]  & \$0.334 & \$0.672$^{*}$ \\
\bottomrule
\multicolumn{8}{@{}p{0.97\linewidth}@{}}{\footnotesize $^{\dagger}$Operating
point: the threshold used for every \sysname number in the paper.
$^{*}$\$/solve from a solo-filled point estimate of the solve count, since
tasks routed to a fixer with no measured with-handoff outcome take that
fixer's solo outcome; rows whose solve count is not measured exactly report
the bracketing interval instead, and no midpoint. Rows above the operating
point send tasks that pass no gate to the most capable fixer (\fxopus); at the
operating point every task passes at least one gate, so the fall-through never
fires. Solves are out of $n{=}266$.
All rows use the board's cost convention: measured fixer API spend plus a
pinned \$5.13 of searcher GPU and sandbox infrastructure, amortized over 266
tasks.}
\end{tabular}%
\endgroup

\end{table}



% ============================================================================
% C — Verify-then-Strip Outcome Classes
% ============================================================================
\section{Verify-then-Strip Outcome Classes}
\label{apx:verify-strip}

\S\ref{sec:results} reports the headline guard statistics; the full
failure-class census behind those numbers appears in
Table~\ref{tab:apx-verify-strip}, broken down by outcome category and by
handoff type.

\begin{table}[H]
\centering
\footnotesize
\caption{\textbf{What the verify-then-strip sandbox found in 266 handoffs.}
\emph{Top:} outcome class of each reproduction claim, replayed before any
fixer saw it ($50$ genuine, $174$ false and stripped).  \emph{Bottom:}
split by handoff type; spontaneous handoffs are more than twice as likely
to be genuine as forced ones (never pooled).  Forced-handoff percentages
use the $43$ of $60$ that claim a reproduction as their denominator.
The label-run row shows the same measurement on an earlier, easier task
set.}
\label{tab:apx-verify-strip}
% ===========================================================================
% tab-apx-verify-strip.tex — APPENDIX. Verify-then-strip outcomes over all
% 266 search-phase handoffs, plus the spontaneous/forced split and the
% label-run baseline.
% BODY ONLY (tabular). Wrapper/caption/label live in the appendix file.
%
% PROVENANCE — every cell.
%   Source: internal records (frozen; run over all
%   266 handoffs with verify_strip_reconstruct.py against live Pro-image
%   sandboxes). Verified verbatim:
%     "Total handoffs processed: 266"
%     "Claiming reproduced: true: 249/266"
%     "Genuine reproductions (real_failure among claim=true): 50/249 (20%)"
%     "STRIPPED as false (passed_stripped among claim=true): 174/249 (70%)"
%   Failure-class table (all 266), verbatim:
%     passed_stripped 174 | real_failure 50 | claim_not_true 16 |
%     import_error 12 | nonzero_other 9 | missing_file 3 | no_repro 1 |
%     no_command 1     (sums to 266 — checked)
%   Spontaneous vs forced table (NEVER pooled), verbatim:
%     spontaneous n=206, all 206 claiming, genuine 46/206 (22%),
%       stripped 140/206 (68%)
%     forced n=60, 43 claiming, genuine 4/43 (9%), stripped 34/43 (79%)
%     NOTE the denominators differ by arm and the source is explicit about
%     it: the spontaneous percentages are over all 206 handoffs (all of which
%     claim), the forced percentages are over the 43 that claim, not over 60.
%     The table below reproduces the source's own denominators and the
%     caption states them.
%   Label-run baseline row, verbatim from the same file's comparison table:
%     "label-run baseline ... 91 claiming | 29 genuine | 32% | 51 false | 56%"
%   Cross-check: internal records restates the
%   same numbers ("249/266 handoffs claimed reproduced. Genuine: 50 (20%).
%   Demonstrably false: 174 (70%) ... Forced handoffs: 9% genuine vs
%   spontaneous 22% ... Worse than the label run's 32%/56%").
%
%   NEVER-POOL RULE (harness rule):
%   spontaneous and forced handoffs are never pooled in any claim.
% ===========================================================================
\begin{tabular}{@{}lr@{}}
\toprule
Outcome of the replayed reproduction & $n$ \\
\midrule
Passed at base (claim stripped)     & 174 \\
Genuinely failed (claim kept)       & \phantom{0}50 \\
Claim not true                      & \phantom{0}16 \\
Import error                        & \phantom{0}12 \\
Non-zero exit, other                & \phantom{00}9 \\
Missing file                        & \phantom{00}3 \\
No reproduction block               & \phantom{00}1 \\
No command                          & \phantom{00}1 \\
\midrule
Total                               & 266 \\
\bottomrule
\end{tabular}

\vspace{4pt}

\begin{tabular}{@{}lrrr@{}}
\toprule
Handoff set & Claiming & Genuine & Stripped \\
\midrule
Spontaneous ($n{=}206$) & 206 & 22\% & 68\% \\
Forced ($n{=}60$)       & \phantom{0}43 & \phantom{0}9\% & 79\% \\
\midrule
All Pro handoffs ($n{=}266$) & 249 & 20\% & 70\% \\
Label-run baseline           & \phantom{0}91 & 32\% & 56\% \\
\bottomrule
\end{tabular}

\end{table}



% ============================================================================
% D — Localization by Repository
% ============================================================================
\section{Localization by Repository}
\label{apx:localization}

The per-task mean localization quality reported in Table~\ref{tab:localization}
pools three repositories that differ substantially in how many files their
gold patches touch.  Table~\ref{tab:apx-localization-repo} disaggregates by
repository and provides the solved-versus-unsolved cross-tab behind the
main-text numbers.

\begin{table}[H]
\centering
\footnotesize
\caption{\textbf{Localization by repository, and against outcome.}
\emph{Top:} \modelname's file-localization quality per repository; spread
tracks the number of gold files per task.  \emph{Bottom:} quality split
by whether the system solved the task.  Better localization correlates
weakly with solving ($+0.024$ recall), so localization is not a gate on
the outcome.}
\label{tab:apx-localization-repo}
% ===========================================================================
% tab-apx-localization-repo.tex — APPENDIX. Per-repository localization
% quality, plus the localization-vs-outcome cross-tab.
% BODY ONLY (tabular). Wrapper/caption/label live in the appendix file.
%
% PROVENANCE — every cell.
%   Source: internal records, per-repo section, verified verbatim:
%     ansible      (n=96): recall 0.472 | precision 0.846 | all-gold 10.4%
%                          | median files named 1.0 | mean gold files 3.80
%     openlibrary  (n=91): recall 0.574 | precision 0.766 | all-gold 28.6%
%                          | median 2.0 | mean gold 3.58
%     qutebrowser  (n=79): recall 0.671 | precision 0.853 | all-gold 38.0%
%                          | median 2.0 | mean gold 2.84
%   Localization vs outcome, derived from data/receipts/ (routing_scenarios.json
%   decisions x outcomes_per_task.json system_union_pairs for the solved set,
%   localization.json per-task rows for the metrics):
%     solved   (n=159): recall 0.575 | precision 0.830 | all-gold 26.4%
%     unsolved (n=107): recall 0.552 | precision 0.808 | all-gold 22.4%
%     The 159/107 split is the routed cell: each task scored on its routed
%     fixer's with-handoff episode, 159/266 solved.
%   Overall (n=266): median files named 2.0, mean files named 2.08,
%     mean gold files 3.44, recall 0.566, precision 0.821, all-gold 24.8%.
%   Repo census n=96/91/79 = the Python-266 census.
%
%   RULE APPLIED: the cross-tab stays
%   on the headline routed cell and is DESCRIPTIVE, not claim-bearing. The
%   caption says so.
%   METHOD STAMP: a single sampled draw per task
%   at temperature 0.9, no majority vote and no best-of-n; gold files are the
%   b-side `diff --git` paths of the gold patch, the same convention the
%   label-run pipeline used.
%   A spontaneous-only cut of the cross-tab removes the handoff-type confound
%   and shows the same modest direction (recall 0.592 solved, n=130 / 0.575
%   unsolved, n=76); not given its own row.
% ===========================================================================
\begin{tabular}{@{}lrrrr@{}}
\toprule
Repository & Recall & Prec. & All-gold & Gold files \\
\midrule
ansible ($n{=}96$)     & 0.472 & 0.846 & 10.4\% & 3.80 \\
openlibrary ($n{=}91$) & 0.574 & 0.766 & 28.6\% & 3.58 \\
qutebrowser ($n{=}79$) & 0.671 & 0.853 & 38.0\% & 2.84 \\
\midrule
All ($n{=}266$)        & 0.566 & 0.821 & 24.8\% & 3.44 \\
\bottomrule
\end{tabular}

\vspace{4pt}

\begin{tabular}{@{}lrrr@{}}
\toprule
Outcome & Recall & Prec. & All-gold \\
\midrule
Solved ($n{=}159$)   & 0.575 & 0.830 & 26.4\% \\
Unsolved ($n{=}107$) & 0.552 & 0.808 & 22.4\% \\
\bottomrule
\end{tabular}

\end{table}



% ============================================================================
% E — Capped-Tier Statistics
% ============================================================================
\section{Capped-Tier Statistics}
\label{apx:cap-stats}

The protocol findings in \S\ref{sec:results} quote selected cap-binding
rates and auto-submit rescue counts.  Table~\ref{tab:apx-cap-stats} gives
the complete picture for all three fixer arms.

\begin{table}[H]
\centering
\footnotesize
\caption{\textbf{How often the evaluation caps bound, and what
auto-submission was worth.}  All arms ran the official capped tier ($50$
turns, \$2.00 per attempt).  The turn cap is the binding constraint; the
cost cap bound only five times.  Auto-submission rescued $24$ of
\fxopus's solves and $33$ of \fxgpt's; without it, both anchors' rates
would fall roughly $9$--$13$ points.  The gold-patch control passes on
$265$ of $266$ tasks.}
\label{tab:apx-cap-stats}
% ===========================================================================
% tab-apx-cap-stats.tex — APPENDIX. What the capped evaluation tier actually
% did: turn-cap and cost-cap incidence, and the value of auto-submission.
% BODY ONLY (tabular). Wrapper/caption/label live in the appendix file.
%
% PROVENANCE — every cell.
%   Source: internal records — per-arm cap behaviour at full scale, plus the
%   addendum covering the cheap-fixer arm. Verified verbatim:
%     gpt-5.2   : "Hit 50-turn cap 97 (36.5%)" | "Hit $2 cost cap 0"
%                 | "Autosubmitted at cap 94" | "Solves rescued by
%                 autosubmit 33" | resolved 139/266 | max attempt $1.5948
%     opus-4.6  : "Hit 50-turn cap 46 (17.3%)" | "Hit $2 cost cap 5"
%                 | "Autosubmitted at cap 48" | "Solves rescued by
%                 autosubmit 24" | resolved 158/266 | max attempt $2.0889
%     kimi-k2.5 : ADDENDUM table, "Caps | 128 step-cap (48%) · 0 cost-cap ·
%                 126 autosubmitted"; resolved 149/266. The addendum does NOT
%                 report a rescued-solve count for this arm, so the table
%                 leaves that cell blank rather than inferring one.
%   Verbatim: "on the full census the cap binds 5 times on opus, and the
%     max attempt reaches $2.0889"; "gpt never hits the cost cap (max $1.59)
%     — its constraint is purely the turn limit, which it hits more than
%     twice as often as opus (36.5% vs 17.3%)".
%   Verbatim on the value of auto-submission: "It rescued 33 gpt solves
%     and 24 opus solves. Without it the headline would read gpt 106/266
%     (39.8%) and opus 134/266 (50.4%)". That is the 9-13pp figure quoted in
%     internal records ("Auto-submit at cap is worth 9-13pp: rescued 24/158 opus, 33/139 gpt solves"): 52.26-39.85 =
%     12.4pp for gpt and 59.40-50.38 = 9.0pp for opus. Both arithmetic
%     checks reproduce the source's range.
%   Gold-control census, verbatim: "Full gold census — 265/266 pass (ansible
%     95/96, openlibrary 91/91, qutebrowser 79/79)"; the single failure is a
%     dataset artifact (a truncated fail-to-pass test id), "unwinnable for
%     every arm equally; kept in all denominators; never a model failure".
%
%   PROTOCOL: capped tier = 50 turns, SILENT (no budget sentence in any
%   prompt), $2.00 per-attempt cost limit, working diff auto-submitted at the
%   cap.
% ===========================================================================
\begin{tabular}{@{}lrrr@{}}
\toprule
 & \fxgpt & \fxopus & \fxkimi \\
\midrule
Hit 50-turn cap        & 36.5\% & 17.3\% & 48\% \\
Hit \$2 cost cap       & 0 & 5 & 0 \\
Max attempt cost       & \$1.59 & \$2.09 & --- \\
Auto-submitted at cap  & 94 & 48 & 126 \\
Solves rescued         & 33 & 24 & --- \\
Solves                 & 139 & 158 & 149 \\
Rate without rescue    & 39.8\% & 50.4\% & --- \\
\bottomrule
\end{tabular}

\end{table}



% ============================================================================
% F — Router Design Space, Full Grid
% ============================================================================
\section{Router Design Space, Full Grid}
\label{apx:design-space}

The calibration analysis in \S\ref{sec:calibration} summarizes the four
feature variants under the logistic-regression scorer.
Table~\ref{tab:apx-design-space} extends that comparison to include the MLP
head and provides the outcome-separation (AUC) measurements the cost numbers
rest on.

\begin{table}[H]
\centering
\footnotesize
\caption{\textbf{The complete router design space.}
\emph{Top:} held-out cost saving for every feature-set $\times$ scorer
combination, all pinned to solve rate $.606$.  Adding hidden states helps
under both scorers; handoff-text features collapse savings to
$8$--$9\%$.  The MLP halves every variant's saving.  At $N{=}99$ the
best variant is only $3/5$ fold-stable.  \emph{Bottom:} the hidden
state separates solved from failed tasks at AUC~$.600$, while
handoff-text embedding sits at chance.}
\label{tab:apx-design-space}
% ===========================================================================
% tab-apx-design-space.tex — APPENDIX. The full router design space: four
% feature variants x two head families, plus the separation (AUC) result the
% cost numbers rest on.
% BODY ONLY (tabular). Wrapper/caption/label live in the appendix file.
%
% PROVENANCE — every cell.
%   Cost-saving rows. Source: internal records — the definitive 4-row
%   comparison on the combined deployable router (LR head) and the head
%   ablation, LR vs the project's tiny MLP on the same 4 rows. Both tables
%   are held-out, fold-wise, N=99, every variant pinned to the SAME solve
%   rate .606, anchor = always-best gpt at $0.406/task. Verified verbatim:
%     variant A (task text only)                   LR +30.5% ($0.2824)
%                                                 MLP +15.4% ($0.3435)
%     variant B (+ handoff text)                   LR  +8.0% ($0.3734)
%                                                 MLP  +5.3% ($0.3845)
%     variant C (+ hidden state, deployed)         LR +34.3% ($0.2668)
%                                                 MLP +18.7% ($0.3299)
%     variant D (+ handoff text + hidden state)    LR  +9.4% ($0.3677)
%                                                 MLP  +4.2% ($0.3890)
%   Fold stability vs anchor: A 3/5, B 1/5,
%     C 3/5, D 4/5. Brier: 0.2698 / 0.2537 / 0.2486 / 0.2505.
%   Verbatim on the head comparison: the MLP "halves every row's saving
%     (C drops 34.3%->18.7%, A drops 30.5%->15.4%)", is "visibly
%     overconfident" (23% of predictions at the 0/1 extremes, reliability gap
%     0.42 on the text-only variant), and "the row ranking is head-invariant
%     where it matters". Recipe pin: "ship LogisticRegression (uncentered)".
%     Same verdict in internal records.
%   Cross-check: the two headline savings, -34.3% (blend, Brier 0.2486,
%     3/5 folds) and -30.5% (text-only, Brier 0.2698, 3/5 folds), reproduce
%     at the same matched solve 0.6061 in internal records.
%
%   Separation rows. Source: internal records,
%   pooled AUC column, verified verbatim:
%     pre-decode layer -4 hidden state : AUC .600, per-fold
%       [.597, .532, .655, .639, .579], AUC-stable 4/5
%     handoff text (the label-run baseline) : AUC .510, per-fold
%       [.436, .439, .752, .424, .501] — "at chance and unstable"
%     issue text : AUC .561
%   The "stable" column is only defined for the two CANDIDATES the source
%   fold-tests against the baseline (the hidden state and the hybrid, both
%   4/5). No fold-stability count exists for the baseline rows themselves, so
%   those cells are dashed rather than filled in.
%   The same section states the honest bound the caption carries: this is an
%   AUC-separation result, and on routing accuracy the hidden state is -3pp
%   and only 3/5 fold-stable — the separation does not convert to picking a
%   better fixer, it converts to pricing the gate (mechanism read:
%   cross-fixer prediction correlation approx 0.78, "a calibrated difficulty
%   gauge, not a per-fixer selector").
%
%   SANITIZED NAMING: the internal documents label these four variants with
%   row letters; here they are named by what they contain (text-only,
%   +handoff-text, text+hidden [deployed], all three) and the letters are
%   kept only as compact labels matching the design-space figure.
%
%   NOT BANKABLE (source's own words): "even its best row holds at only
%   3/5 fold-stability — a real ~4pp improvement that is still not bankable
%   at N=99". The caption must not overstate this.
% ===========================================================================
\begin{tabular}{@{}lrrrr@{}}
\toprule
Router features & \multicolumn{2}{c}{Saving} & Folds & Brier \\
\cmidrule(lr){2-3}
 & LR & MLP & & \\
\midrule
A\quad text only                   & 30.5\% & 15.4\% & 3/5 & 0.270 \\
B\quad $+$ handoff text            & \phantom{0}8.0\% & \phantom{0}5.3\% & 1/5 & 0.254 \\
C\quad $+$ hidden state \emph{(dep.)} & \textbf{34.3\%} & 18.7\% & 3/5 & \textbf{0.249} \\
D\quad $+$ both                    & \phantom{0}9.4\% & \phantom{0}4.2\% & 4/5 & 0.251 \\
\bottomrule
\end{tabular}

\vspace{4pt}

\begin{tabular}{@{}lrr@{}}
\toprule
Routing representation & Pooled AUC & Stable \\
\midrule
Hidden state (pre-decode, layer $-4$) & \textbf{0.600} & 4/5 \\
Task text                             & 0.561 & --- \\
Handoff text (prereg.\ baseline)      & 0.510 & --- \\
\bottomrule
\end{tabular}

\end{table}



% ============================================================================
% G — Redistribution, Exact Numbers
% ============================================================================
\section{Redistribution, Exact Numbers}
\label{apx:redistribution}

Figure~\ref{fig:redistribution} in the main text visualizes the paired
handoff ablation; Table~\ref{tab:apx-redistribution} gives the exact per-fixer
rates, deltas, confidence intervals, discordant-pair counts, and $p$-values
behind that figure.

\begin{table}[H]
\centering
\footnotesize
\caption{\textbf{The handoff pattern is redistributive, not additive.}
Exact per-fixer rates, deltas, and $p$-values behind
Figure~\ref{fig:redistribution} ($99$ paired tasks, $396$ attempts).
The three weaker fixers gain and the strongest loses; pooled effect is
$+1.8$\,pp with CI including zero.  No fixer reaches $p<0.05$; every
row is directional only.  $b$~counts rescued tasks, $c$~counts broken
ones.}
\label{tab:apx-redistribution}
% ===========================================================================
% tab-apx-redistribution.tex — APPENDIX. The controlled paired handoff
% ablation: every fixer run twice on the same tasks, with and without a
% handoff.
% BODY ONLY (tabular). Wrapper/caption/label live in the appendix file.
%
% PROVENANCE — every cell.
%   Source: internal records table, verified
%   verbatim (paired, same task / same fixer, 99 tasks; Delta = handoff minus
%   solo solve rate; exact McNemar = two-sided binomial on the discordant
%   pairs; b = solo-fail -> handoff-solve, c = solo-solve -> handoff-fail):
%     claude : 48.5 -> 53.5 | +5.1 [ 0.0, 10.1] | b 6  | c 1 | p 0.125
%     kimi   : 52.5 -> 56.6 | +4.0 [-4.0, 12.1] | b 11 | c 7 | p 0.481
%     flash  : 55.6 -> 57.6 | +2.0 [-5.1,  9.1] | b 7  | c 5 | p 0.774
%     gpt    : 60.6 -> 56.6 | -4.0 [-11.1, 3.0] | b 5  | c 9 | p 0.424
%     pooled : 54.3 -> 56.1 | +1.8 [-1.0,  4.5] | b 29 | c 22 | p 0.401
%   Scope stamp (binding, from the same file's preamble): N=99 paired /
%   100 solo, Python-only, four fixers, 396 handoff-arm attempts; "N=99-100
%   underpowers every per-fixer significance test — read the CIs, not the
%   point estimates".
%   Verdict line, verbatim: "Not one fixer clears p<0.05 ... every effect is
%   DIRECTIONAL-ONLY", and "gpt's harm (-4pp, 9 regressions vs 5 flips) is
%   reported with identical prominence".
%   Cross-check: internal records restate the same five deltas and the same honest sentence ("handoffs
%   redistribute ability toward weaker/cheaper fixers at no cost change; they
%   do not beat the best solo model at N=99"), and cite this paired set as "the clean controlled citation".
%
%   ORDERING NOTE: rows are printed in ascending solo solve rate, so the
%   redistribution pattern (the weaker the solo fixer, the larger the lift)
%   is legible without the reader doing arithmetic.
% ===========================================================================
\begin{tabular}{@{}lrrrrr@{}}
\toprule
Fixer & Solo & Handoff & $\Delta$ (95\% CI) & $b/c$ & $p$ \\
\midrule
\fxopus & 48.5 & 53.5 & $+5.1$ [$0.0$, $10.1$]  & 6/1  & 0.125 \\
\fxkimi   & 52.5 & 56.6 & $+4.0$ [$-4.0$, $12.1$] & 11/7 & 0.481 \\
\fxflash  & 55.6 & 57.6 & $+2.0$ [$-5.1$, $9.1$]  & 7/5  & 0.774 \\
\fxgpt    & 60.6 & 56.6 & $-4.0$ [$-11.1$, $3.0$] & 5/9  & 0.424 \\
\midrule
pooled & 54.3 & 56.1 & $+1.8$ [$-1.0$, $4.5$]  & 29/22 & 0.401 \\
\bottomrule
\end{tabular}

\end{table}



% ============================================================================
% H — Benchmark-Versus-Fresh Calibration Receipts
% ============================================================================
\section{Benchmark-Versus-Fresh Calibration Receipts}
\label{apx:bench-fresh}

The router's r\'esum\'es are built from public per-task outcomes, so the
router inherits whatever biases those outcomes carry.  Buying fresh labels was
not a design preference but a measured necessity.
Table~\ref{tab:apx-bench-fresh} records the calibration receipts that
motivated that decision: every fixer's public rate compared with its
freshly measured rate under a single controlled harness.

\begin{table}[H]
\centering
\footnotesize
\caption{\textbf{Why we bought fresh labels: public rates do not transfer
to unseen tasks.}  Each fixer's public multilingual rate versus its freshly
measured rate on $100$ post-cutoff Python tasks under one harness.
Every fixer drops; the ordering does not survive.  Rates use the full
$100$-task solo denominator; the paired analysis in
Appendix~\ref{apx:redistribution} uses the $99$ tasks with both arms.
These are calibration receipts, not a benchmark claim.}
\label{tab:apx-bench-fresh}
% ===========================================================================
% tab-apx-bench-fresh.tex — APPENDIX. Calibration receipts: public benchmark
% rates against freshly measured rates for the same four fixers.
% BODY ONLY (tabular). Wrapper/caption/label live in the appendix file.
%
% PROVENANCE — every cell.
%   Source: internal records table, verified
%   verbatim (phase-2 solo data, 100 fresh post-cutoff Python tasks, against
%   the 9-language public Multilingual leaderboard):
%     gpt    : bench .667 (#3)  | fresh .610 [.510, .700] | gap  -5.7
%     flash  : bench .727 (#1)  | fresh .550 [.450, .650] | gap -17.7
%     kimi   : bench .673 (#2)  | fresh .520 [.420, .610] | gap -15.3
%     claude : bench .720 (#1=) | fresh .480 [.380, .570] | gap -24.0
%   Fresh ranks are the source's own ordering line: "fresh-solo order: gpt >
%   flash > kimi > claude", i.e. gpt #1, flash #2, kimi #3, claude #4.
%   Source's read, verbatim: "Every fixer drops; the two benchmark leaders
%   (flash .727, claude .720) fall hardest, and the leaderboard ordering does
%   not survive contact with unseen post-cutoff tasks — gpt was #3 on the
%   bench and is #1 fresh".
%   NOTE ON THE FRESH RATE FOR gpt: the benchmark-vs-fresh comparison gives
%   .610 while the paired analysis gives .606 for the same fixer on the
%   99-task paired subset. Both are in the same source and are not in
%   conflict — different denominators. This table prints .610, the
%   benchmark-vs-fresh figure.
%
%   FRAMING RULE (BINDING): "benchmark-vs-fresh gap =
%   calibration receipts only, never a claim". The caption below is written
%   to that rule: this table exists to explain why fresh labels were bought,
%   NOT to argue that public benchmarks are inflated. Do not re-word it into
%   a claim about benchmark integrity. The absolute levels are also
%   language-caveated (Python-only fresh set vs a 9-language board) — the
%   source states that only the cross-fixer ORDERING is the comparable part.
% ===========================================================================
\begin{tabular}{@{}lrrrr@{}}
\toprule
Fixer & Public & Fresh (95\% CI) & Gap & Rank \\
\midrule
\fxflash & 0.727 & 0.550 [.450, .650] & $-17.7$ & \#1 $\rightarrow$ \#2 \\
\fxopus  & 0.720 & 0.480 [.380, .570] & $-24.0$ & \#1 $\rightarrow$ \#4 \\
\fxkimi  & 0.673 & 0.520 [.420, .610] & $-15.3$ & \#2 $\rightarrow$ \#3 \\
\fxgpt   & 0.667 & 0.610 [.510, .700] & $-5.7$ & \#3 $\rightarrow$ \#1 \\
\bottomrule
\end{tabular}

\end{table}



% ============================================================================
% I — Handoff Effect on Router-Selected Benchmark Subsets
% ============================================================================
\section{Handoff Effect on Router-Selected Benchmark Subset}
\label{apx:pro-handoff}

The redistribution pattern observed under controlled conditions
(\S\ref{sec:calib-redist}) also appears in the benchmark evaluation, where
the tasks the router sends to \fxkimi were available both with and without a
handoff.  Table~\ref{tab:apx-pro-handoff} reports that comparison.  The subset
was selected by the router rather than randomly assigned; it remains measured,
uncontrolled corroboration rather than a controlled experiment.

\begin{table}[H]
\centering
\footnotesize
\caption{\textbf{The handoff effect on the deployment path.}
The tasks routed to \fxkimi were also run without a handoff, giving a
same-task comparison: the workhorse gains $3.8$\,pp on $263$ tasks.
$b/c$ counts tasks solved only with the handoff versus only solo; the
exact McNemar test on those discordant pairs gives $p=0.245$.  The
router selected \emph{which} tasks enter the subset, not the
within-task contrast, so the pairing is intact, but the subset is not
randomly assigned and this table remains corroboration, not a
controlled measurement.  The three tasks routed elsewhere have no solo
arm for their fixer and are omitted.}
\label{tab:apx-pro-handoff}
% ===========================================================================
% tab-apx-pro-handoff.tex — APPENDIX. Deployment-path corroboration of the
% handoff effect, on the router-selected subset of the Pro evaluation.
% BODY ONLY (tabular). Wrapper/caption/label live in the appendix file.
%
% PROVENANCE — every cell. Derived from data/receipts/:
%   routing_scenarios.json `decisions` (theta = 0.30) gives the routed subset:
%     263 tasks -> kimi-k2-5, 3 -> gemini-3-flash.
%   outcomes_per_task.json:
%     with handoff = system_union_pairs[task]["kimi-k2-5"].resolved
%       -> 158/263 = 60.1%
%     solo, same tasks = solo[task]["kimi"].resolved
%       -> 148/263 = 56.3%
%     delta = +3.8pp
%     discordant pairs: handoff-only b=35, solo-only c=25 (both 123,
%       neither 80); exact two-sided McNemar on b+c=60 -> p = 0.245
%   Whole-benchmark anchors from the same file: Kimi K2.5 with handoff
%   159/266, solo 149/266.
%
%   WHY ONE ROW: the 3 tasks routed to Gemini 3 Flash have no solo Flash arm,
%   so they cannot form a paired subset; the caption's footnote says so.
%   The subset is router-selected, not randomly assigned, so the paired
%   label-run comparison stays the controlled citation and this table is
%   corroboration only; the McNemar p is printed with that caveat in the
%   caption.
% ===========================================================================
\begin{tabular}{@{}lrrrrrr@{}}
\toprule
Fixer & Routed $n$ & Handoff & Solo & $\Delta$ & $b/c$ & $p$ \\
\midrule
\fxkimi & 263 & 60.1\% & 56.3\% & $+3.8$ & 35/25 & 0.245 \\
\bottomrule
\end{tabular}

\end{table}



% ============================================================================
% J — Nine-Language Localization Detail
% ============================================================================
\section{Nine-Language Localization Detail}
\label{apx:languages}

Figure~\ref{fig:language-transfer} in the main text shows the per-language
$F_1$ bars; Table~\ref{tab:apx-languages} gives the full precision, recall,
$F_1$, and all-gold counts behind that figure, split by trained and
never-trained language pools.

\begin{table}[H]
\centering
\footnotesize
\caption{\textbf{Localization transfers to languages \modelname never saw
in training.}  Per-language precision, recall, and $F_1$ of spontaneous
handoffs against gold patch files.  The six never-trained languages
outperform the three trained ones ($F_1$~$0.630$ vs.\ $0.455$).
JavaScript and TypeScript are genuine soft spots; the TypeScript and C++
cells rest on $12$ assigned tasks each, with $7$ and $8$ spontaneous
handoffs analyzed ($\dagger$).  Multi-file recall is weak across
all languages ($0.26$--$0.34$).  Forced handoffs excluded; searcher
alone, one sampled draw.}
\label{tab:apx-languages}
% ===========================================================================
% tab-apx-languages.tex — APPENDIX. Per-language localization quality of
% spontaneous handoffs, across three trained and six never-trained languages.
% BODY ONLY (tabular). Wrapper/caption/label live in the appendix file.
%
% PROVENANCE — every cell.
%   Source: internal records, Table 4a (per-language)
%   and Table 4b (trained vs never-trained aggregate). Every stamp on those
%   tables applies here: single draw, temperature 0.9, solo searcher with no
%   fixer in the loop, gold-patch-files metric only, SWE-bench Multilingual;
%   spontaneous and forced are NEVER pooled and only the spontaneous arm is
%   printed below.
%   Table 4a SPONT columns, verified verbatim (P / R / F1, all-hit):
%     Go         trained  n=42  spont 27 : 0.605 / 0.542 / 0.546  13/27
%     JavaScript trained  n=31  spont 26 : 0.500 / 0.365 / 0.389   7/26
%     TypeScript trained  n=12  spont  7 : 0.333 / 0.357 / 0.343   2/7
%     Java       never    n=43  spont 27 : 0.722 / 0.574 / 0.612  11/27
%     Ruby       never    n=44  spont 36 : 0.833 / 0.661 / 0.701  18/36
%     Rust       never    n=43  spont 20 : 0.800 / 0.626 / 0.667  10/20
%     PHP        never    n=42  spont 33 : 0.697 / 0.632 / 0.630  17/33
%     C          never    n=30  spont 18 : 0.546 / 0.435 / 0.435   6/18
%     C++        never    n=12  spont  8 : 0.688 / 0.813 / 0.708   6/8
%   Table 4b, spontaneous rows, verified verbatim:
%     trained (Go/JS/TS)              n=60  : P 0.528 R 0.444 F1 0.455
%                                             all-hit 22/60 (37%)
%     never-trained (6 languages)     n=142 : P 0.731 R 0.613 F1 0.630
%                                             all-hit 68/142 (48%)
%   Difficulty-matched control (source's difficulty-matched cut,
%   spontaneous, single-gold-file tasks — the strata that isolate "did you
%   name the one right file"): trained 42 tasks, recall 0.524; never-trained
%   86 tasks, recall 0.791. (Multi-gold strata, also in the source: 0.257
%   trained vs 0.339 never-trained — weak everywhere, quoted in the caption
%   as the honest bound rather than tabulated.)
%   INDICATIVE-ONLY MARK: the source stamps TypeScript and C++ (n=12 episodes
%   each, 7 and 8 spontaneous) as indicative-only. Both rows carry a dagger
%   here and the caption repeats the warning.
%   Trained pool = Go, JavaScript, TypeScript (Python is absent from this
%   task pool by construction); never-trained = Java, Ruby, Rust, PHP, C, C++
%   (source's method note).
%   Licensed phrasing: "transfers to six
%   never-trained languages with localization quality equal or better,
%   difficulty-matched; JS/TS is the trained-pool weak spot." NOT licensed:
%   any multi-file or reproduction-quality claim. The caption respects both.
% ===========================================================================
\begin{tabular}{@{}lrrrrr@{}}
\toprule
Language & $n$ & P & R & $F_1$ & All-gold \\
\midrule
Go                  & 27 & 0.605 & 0.542 & 0.546 & 13/27 \\
JavaScript          & 26 & 0.500 & 0.365 & 0.389 & \phantom{0}7/26 \\
TypeScript$^{\dagger}$ & \phantom{0}7 & 0.333 & 0.357 & 0.343 & \phantom{0}2/7 \\
\addlinespace
\emph{trained pool} & 60 & 0.528 & 0.444 & 0.455 & 22/60 \\
\midrule
Java                & 27 & 0.722 & 0.574 & 0.612 & 11/27 \\
Ruby                & 36 & 0.833 & 0.661 & 0.701 & 18/36 \\
Rust                & 20 & 0.800 & 0.626 & 0.667 & 10/20 \\
PHP                 & 33 & 0.697 & 0.632 & 0.630 & 17/33 \\
C                   & 18 & 0.546 & 0.435 & 0.435 & \phantom{0}6/18 \\
C++$^{\dagger}$     & \phantom{0}8 & 0.688 & 0.813 & 0.708 & \phantom{0}6/8 \\
\addlinespace
\emph{never-trained} & 142 & 0.731 & 0.613 & \textbf{0.630} & 68/142 \\
\bottomrule
\end{tabular}

\end{table}



% ============================================================================
% K — Vault Decoding Comparison
% ============================================================================
\section{Vault Decoding Comparison}
\label{apx:vault}

The decoding finding reported in \S\ref{sec:decoding} is demonstrated on
a 100-task dial subset; Table~\ref{tab:apx-vault} gives the full A/B
comparison on the 450-task held-out vault that established the finding.

\begin{table}[H]
\centering
\footnotesize
\caption{\textbf{Sampled decoding is what makes the searcher commit.}
The same checkpoint on $450$ held-out issues, greedy versus temperature
$0.9$.  Find rate nearly triples, from higher commitment
($72\%$ vs.\ $21\%$) against an 18\% drop in per-handoff recall.  The
greedy arm lost $24$ episodes to infrastructure timeouts; the
$2.65\times$ figure adjusts for this.  The vault is excluded from all
training data.}
\label{tab:apx-vault}
% ===========================================================================
% tab-apx-vault.tex — APPENDIX. Sampled versus greedy decoding for the
% searcher, on the 450-task held-out vault.
% BODY ONLY (tabular). Wrapper/caption/label live in the appendix file.
%
% PROVENANCE — every cell.
%   Source: internal records (the 450-vault
%   milestone comparison, run A = temperature 0.9, run B = greedy, same
%   checkpoint, same 450 tasks). Verified verbatim from the metric table:
%     n                450 / 450
%     find_rate        0.3058 / 0.1104   (A/B 2.77x raw)
%     emission         0.7178 / 0.2133
%     handoffs         323 / 96
%     recall/handoff   0.4261 / 0.5173
%     precision        0.6428 / 0.7257
%     files/handoff    1.93 / 2.05
%     wrong-commit     0.195 / 0.1354
%     all-files-hit    71 / 27
%     malformed        0 / 0
%   Error bars, verbatim: "find_rate = 0.3058  95% CI [0.2731, 0.3371]
%     half-width +-0.032" (bootstrap 1000x over the 450 tasks, seed 0).
%   Timeout sensitivity, verbatim: "RAW A=0.3058 B=0.1104 ratio 2.77x /
%     MATCHED A=0.3091 B=0.1166 ratio 2.65x (n=426)" — the greedy arm lost 24
%     episodes to infrastructure wall-timeouts, which are dropped from BOTH
%     arms for the matched figure.
%   THE QUOTABLE MULTIPLIER IS 2.65x, NEVER 5.7x. Two separate internal
%     retractions bind this:
%     (a) internal records: "the 0.0631
%         comparator was ONE CONTAMINATED RUN ... DO NOT USE; 5.7x claim
%         RETRACTED";
%     (b) internal records: the honest multiplier is 2.65x (timeout-matched,
%         dropping 24 greedy-arm timeouts from both arms; raw 2.77x is not
%         quoted, since greedy timeouts flatter us).
%     The raw 2.77x is deliberately NOT the headline for the same reason.
%   Decomposition (quoted in the caption): "commitment
%     3.37x x per-handoff quality 0.82x — the vault charges a real quality
%     give-back". The quality give-back is visible in this table's own
%     recall-per-handoff and precision rows, which both fall.
%   Vault construction: 450 issues, Go 200 / TS 150 / Py 100 over 306 repos,
%     deterministic selection, and the vault ids are a mandatory exclusion
%     list for every training set.
% ===========================================================================
\begin{tabular}{@{}lrr@{}}
\toprule
Metric ($n{=}450$) & Sampled & Greedy \\
\midrule
Find rate            & 0.306 & 0.110 \\
\quad 95\% CI        & [.273, .337] & --- \\
\quad timeout-matched ratio & \multicolumn{2}{r}{$2.65\times$ ($n{=}426$)} \\
\addlinespace
Emission rate        & 0.718 & 0.213 \\
Handoffs emitted     & 323 & \phantom{0}96 \\
\addlinespace
Recall per handoff   & 0.426 & 0.517 \\
Precision            & 0.643 & 0.726 \\
Files per handoff    & 1.93 & 2.05 \\
All-gold handoffs    & \phantom{0}71 & \phantom{0}27 \\
\addlinespace
Malformed            & \phantom{00}0 & \phantom{00}0 \\
\bottomrule
\end{tabular}

\end{table}



% ============================================================================
% L — Training Hyperparameters
% ============================================================================
\section{Training Hyperparameters}
\label{apx:hyperparams}

\S\ref{sec:training} describes the supervised recipe at the level
relevant to the contribution; Table~\ref{tab:apx-hyperparams} records the
full configuration as realized in the actual training run.

\begin{table}[H]
\centering
\footnotesize
\caption{\textbf{Supervised fine-tuning configuration, as run.}
Rank-$64$ LoRA on every linear projection, trained in \code{bf16} for
two epochs on packed $32$k-token blocks at $91.7\%$ fill; loss is on
assistant turns only.  All values are the realized run, not the planned
recipe.  The complete run cost \$86 of GPU time and peaked at
$24.2$\,GB.}
\label{tab:apx-hyperparams}
% ===========================================================================
% tab-apx-hyperparams.tex — APPENDIX. Supervised fine-tuning configuration
% for \modelname, as actually run.
% BODY ONLY (tabular). Wrapper/caption/label live in the appendix file.
%
% PROVENANCE — every cell. Where a planned value and a realized value differ,
% the REALIZED value is printed and the divergence is recorded here.
%   REALIZED source: internal records: "Locked config (measured, no
%   further tuning)", verbatim:
%     "Qwen2.5-Coder-7B-Instruct · LoRA r64/alpha128 all-linear · bf16 ·
%      LR 1e-4 (cosine, 10% floor, warmup 0.03) · 2 epochs · 32k · TRUE
%      packed blocks · mbs 1 · grad-accum 8 · grad-ckpt ON · xformers ·
%      seed 3407 · save-steps 20 (resumable, keep-2) ... 3,182 steps
%      (1,591/epoch)"
%   internal records also records the adapter shape ("adapter: r64/alpha128,
%     use_rslora=F, modules_to_save=None, base=unsloth/Qwen2.5-Coder-7B-
%     Instruct" — i.e. no embedding or lm_head modules trained) and the
%     realized run totals: "Training done: 3182 steps / 2.0 epochs / 760M
%     tokens / wall 31.8h / peak VRAM 24.2GB / ~$86. Final loss 0.114 (step
%     3182)."
%   Per-epoch token count and packing efficiency, internal records: "Run
%     identity": "12,723 blocks / 19,905 ex / 382,417,805 tok / 91.7% fill".
%   PLANNED source: internal records ("bf16 LoRA — r=64, alpha=128, ALL
%     linear projections (q/k/v/o+gate/up/down), no embed/lm_head, dropout 0,
%     LR 1e-4-2e-4 (10x full-FT rule), 2-3 epochs, 32k ctx with sequence
%     packing + gradient checkpointing. NOT QLoRA") and internal records ("micro-batch
%     1x32k packed ..., grad-accum 8 ...; warmup 3%, cosine to 10% of peak;
%     LR 1e-4 ...; 2 epochs (not 3 — behavioral overfitting precedes loss
%     signal)"). Same recipe fixed in internal records.
%
%   DIVERGENCES, all resolved in favour of the realized run, none material:
%     LR      : planned range 1e-4 to 2e-4  -> realized 1e-4 (bottom of range)
%     Epochs  : planned 2-3                 -> realized 2 (internal records already
%               narrowed this to 2 pre-run, with the reason stated)
%     Packing : the plan assumed packing was active; internal records note
%               that the library's packing import failed silently and that
%               TRUE packing was reimplemented in the project's own
%               preprocessing before the paid run. The 91.7% fill figure is
%               the measured fill of that implementation, not an estimate.
%     Attention backend: xformers, not flash-attention — internal records
%               defers the flash-attention migration explicitly ("one-change-
%               at-a-time before paid run").
%   Every planned value has a realized counterpart in internal records.
% ===========================================================================
\begin{tabular}{@{}ll@{}}
\toprule
Setting & Value \\
\midrule
Base model        & Qwen2.5-Coder-7B-Instruct \\
Adaptation        & LoRA, $r{=}64$, $\alpha{=}128$ \\
Target modules    & all linear projections \\
                  & (no embedding, no output head) \\
Dropout           & 0 \\
Precision         & \texttt{bf16} \\
\midrule
Learning rate     & $1{\times}10^{-4}$, cosine to $10\%$ \\
Warmup            & 3\% of steps \\
Epochs            & 2 (3,182 steps) \\
Context           & 32k, packed blocks \\
Micro-batch       & $1{\times}32$k \\
Gradient accum.   & 8 \\
Gradient ckpt.    & on \\
Seed              & 3407 \\
\midrule
Tokens per epoch  & 382.4M (91.7\% fill) \\
Wall clock        & 31.8\,h \\
Peak memory       & 24.2\,GB \\
Final loss        & 0.114 \\
Cost              & \$86 \\
\bottomrule
\end{tabular}

\end{table}

\paragraph{Reinforcement-learning rig.}
We built and validated a GRPO training rig before concluding that RL was
unnecessary (\S\ref{sec:training}).  The rig used a two-GPU topology: one GPU
trained while the other generated rollouts, with kill-and-resume
checkpointing so that preempted runs lost no more than a single episode.
Exact backpropagation through long episodes was verified against ground truth
before any paid run began.  The reward function was validated separately, by
replaying candidate reward assignments against thousands of previously saved
episodes and confirming agreement before committing GPU spend.  50 clean
steps then ran flat, as reported in \S\ref{sec:training}, and the rig was
shelved in favor of the decoding fix.



% ============================================================================
% M — A Verified Handoff, Before and After the Gate
% ============================================================================
\section{A Verified Handoff, Before and After the Gate}
\label{apx:handoff-example}

Figure~\ref{fig:apx-handoff-example} reproduces one real spontaneous handoff
from the SWE-bench Pro evaluation, shown in the exact post-strip form the
fixer received.  This handoff's reproduction claim was replayed in the task's
own sandbox and genuinely failed, so the claim was forwarded intact rather
than deleted.

\begin{figure}[H]
\centering
\footnotesize
% ===========================================================================
% fig-apx-handoff-example.tex — APPENDIX. One real handoff, in the exact form
% the fixer received it.
% BODY ONLY (framed verbatim listing). Wrapper/caption/label live in the
% appendix file.
%
% PROVENANCE — the text below is a real record, not a reconstruction.
%   Source: internal records, the row whose task id
%   begins instance_qutebrowser__qutebrowser-bedc9f7fadf93f83d8dee95feeecb992
%   2b6f063f (SWE-bench Pro, Python-266 census). The text shown is that row's
%   post-strip injection-ready field — i.e. exactly what the routed fixer was
%   given, after the verify-then-strip stage ran. That file is the
%   injection-ready handoff set for the fixer stage.
%
%   WHY THIS ONE (selection criteria, applied in order):
%     (1) SPONTANEOUS — the searcher stopped on its own; the record's handoff
%         type field is "spontaneous". Forced handoffs are never shown as
%         representative.
%     (2) GENUINE REPRODUCTION — the record's failure class is "real_failure",
%         i.e. one of the 50 of 249 claims that actually reproduced when
%         replayed in the task sandbox, so nothing in the reproduction block
%         was stripped. Its verification record reads verified true, stripped
%         false, reason "verified_reproduces", return code 1.
%     (3) SHORT AND CODE-FREE — 1,343 characters, no code dump; it is the
%         fifth-shortest of the 46 spontaneous-and-genuine handoffs, and the
%         shorter ones are less illustrative (single file, one-line notes).
%
%   FIDELITY — two mechanical changes, both disclosed, no content added or
%   removed:
%     (a) LINE WRAPPING: the reproduction block's last line and the notes
%         paragraph are hard-wrapped at 58 columns so the listing fits an
%         IEEE column. Only newlines were inserted; no character was added,
%         deleted or reordered.
%     (b) ONE CHARACTER SUBSTITUTED: the notes contain a U+2011 non-breaking
%         hyphen in "top-level" (an artifact of the open-weights notes
%         writer). It is rendered here as an ordinary hyphen because the
%         document's font encoding has no glyph for it.
%   Nothing is redacted.
%
%   WHAT THE READER SHOULD SEE (for the caption, all checkable against the
%   text): four ranked candidate files; a reproduction command with its
%   observed failure output, which survived the sandbox replay and is
%   therefore forwarded rather than deleted; free-prose notes that record
%   both what was found and what was searched for and not found (the grep
%   that returned nothing) — i.e. dead ends are part of the artifact; and the
%   terminating STOP token that ends the search episode.
%   Handoff schema: internal records (ranked implicated files; failing
%   reproduction with command/reproduced/observed; dead ends and repo notes;
%   XML-style tags, raw unescaped content, <=4KB).
% ===========================================================================
% NOTE: the listing is delimited by rules rather than \fbox{} because a
% verbatim environment cannot be read as a macro argument.
\begin{minipage}{\linewidth}
\rule{\linewidth}{0.6pt}
\scriptsize
\begin{verbatim}
<files>
qutebrowser/utils/utils.py
qutebrowser/utils/qtutils.py
qutebrowser/browser/downloads.py
qutebrowser/mainwindow/tabbedbrowser.py
</files>
<repro>
command: cd /app && python test_issue.py
reproduced: true
observed: OBSERVATION:
SUCCESS: interpolate_color exists in utils
Interpolation at 50%: 128:0:128
Traceback (most recent call last):
  File "/app/test_issue.py", line 26, in <module>
    interp = qtutils.interpolate_color
AttributeError: module 'qutebrowser.utils.qtutils' has no
attribute 'interpolate_color'
</repro>
<notes>
Search surfaced definitions of interpolate_color in
qutebrowser/utils/utils.py (line 260) and references to
utils.interpolate_color in
qutebrowser/browser/downloads.py (line 563) and
qutebrowser/mainwindow/tabbedbrowser.py (lines 866, 883).
No occurrences of interpolate_color were found in
qutebrowser/utils/qtutils.py. Grep for
_get_color_percentage returned nothing. Attempted import
of qutebrowser.utils.qtutils showed only top-level
symbols, confirming the function is absent there. Running
a custom script test_issue.py via cd /app && python
test_issue.py raised AttributeError: module
'qutebrowser.utils.qtutils' has no attribute
'interpolate_color'. No build or test tooling issues
observed; commands executed directly with python. No false
leads beyond missing function in qtutils.
</notes>
STOP
\end{verbatim}
\vspace{-6pt}
\rule{\linewidth}{0.6pt}
\end{minipage}

\caption{\textbf{One real handoff, exactly as the fixer received it.}
A spontaneous handoff from the SWE-bench Pro evaluation, reproduced
verbatim after the verify-then-strip stage.  This reproduction claim was
genuine and forwarded intact; $174$ of $249$ claims were false and
stripped.  Line wrapping in the notes is ours; no content is added or
redacted.}
\label{fig:apx-handoff-example}
\end{figure}



% ============================================================================
% N — Compute and Cost Disclosure
% ============================================================================
% Layout only: start N in a fresh column so its heading, prose and pinned
% table stay together. Without this the heading and prose fall at the foot
% of the previous column and Table XVIII lands alone on the next page.
\newpage
\section{Compute and Cost Disclosure}
\label{apx:spend}

Table~\ref{tab:apx-spend} itemizes the measured spend across the three eras
of this project: evaluation, label collection, and training.
The one-off cost of building the system, comprising the label run
(${\approx}\$370$), the fine-tuning run (${\approx}\$86$), and corpus
note generation (${\approx}\$9$), totals roughly \$465.  Against the
\$0.62 per-task saving over the best solo fixer (\$0.757 API-only
versus \$0.137 all-in), that sunk cost amortizes in roughly 750 tasks.

\begin{table}[H]
\centering
\footnotesize
\caption{\textbf{Measured spend.}
Every figure is a measured value from the run ledgers.  The evaluation
era covers all five arms plus infrastructure; the label-run era covers
fresh per-task outcomes; the training era is the single SFT run.
\modelname's entire evaluation contribution was \$1.13 of GPU time
(${\sim}0.4$ cents per task).  Era totals are not summed.}
\label{tab:apx-spend}
% ===========================================================================
% tab-apx-spend.tex — APPENDIX. Compute and API spend disclosure.
% BODY ONLY (tabular). Wrapper/caption/label live in the appendix file.
%
% PROVENANCE — every line item. Ranges are printed as ranges wherever the
% source gives a range; nothing is rounded into a point estimate.
%
%   EVALUATION ERA — internal records: "Ledger (era,
%   against the $1,200 hard cap)", verified verbatim, item by item:
%     probe (40 attempts)                              19.24
%     gpt solo arm                                    151.64
%     opus solo arm                                   201.25
%     kimi solo arm                                    28.23
%     union331 (system)                                86.49
%     smoke fixers                                      1.32
%     unrecorded (restart-killed partials, disclosed)   ~11
%     RunPod GPU (all pods)                             ~5.9
%     Modal compute (census, sandboxes, drivers,
%       verify-strip)                                   ~40-55
%     era total                                        ~$545-560
%   The "unrecorded" line is the source's own disclosure of restart-killed
%   partial attempts and is reproduced rather than absorbed into another row.
%   Cross-check on the arms: internal records
%   reports the same $151.64 and $201.25, and its addendum the same $28.23,
%   with "Era pinned spend after three arms: $381.12".
%
%   LABEL-RUN ERA — internal records, verbatim:
%     "Total era spend ~= $370 (fixer API $340.03 + Scout pod ~$4 + Modal
%      ~$25) against the $750 ceiling / $400 signed budget", and §Spend:
%     "phase-2 $164.70 · shakedown $11.24 · handoff arm $175.33 -> API
%      $340.03 (episodes cross-checked to the cent vs litellm on 2 providers);
%      + smokes/probe ~$18 · Scout pod ~$4 · Modal ~$25".
%     NOTE: the ~$370 headline and the itemization differ by the ~$18 of
%     smokes/probe (340.03 + 4 + 25 = $369.03 without them). The table prints
%     the source's own ~$370 total and its three named components; the ~$18
%     smoke/probe line is disclosed in this comment and folded into the
%     "approximately" of the headline exactly as the source does.
%
%   TRAINING ERA — internal records: "Training done: 3182 steps /
%     2.0 epochs / 760M tokens / wall 31.8h / peak VRAM 24.2GB / ~$86";
%     the pre-run locked config projected "~31.7 h -> ~$85", and internal
%     records confirm "SFT COMPLETE (3,182 steps, $86, loss 0.114)".
%     Data-side notes generation was ~$9 (run-1 complete: "cost ~$9").
%
%   SEARCHER COST INSIDE THE EVALUATION — internal records: "266
%     search episodes + hidden-state extraction: $1.13 GPU total (~$0.004/task,
%     ~1.3% of system cost/task)". Reported in the caption because it is the
%     number that matters for the economics argument.
%
%   DISCLOSURE RULE (sober disclosure): NO grand-total line beyond what
%   the sources themselves state. Each era carries its own source-stated
%   total; the eras are NOT summed into a single project budget figure, and
%   no narrative is built on the numbers.
% ===========================================================================
\begin{tabular}{@{}lr@{}}
\toprule
Item & \$ \\
\midrule
\multicolumn{2}{@{}l}{\emph{Evaluation era}} \\
Probe (40 attempts)        & 19.24 \\
\fxgpt solo arm           & 151.64 \\
\fxopus solo arm          & 201.25 \\
\fxkimi solo arm         & 28.23 \\
System arm                 & 86.49 \\
Smoke                      & 1.32 \\
Unrecorded partials        & $\approx$11 \\
GPU (all pods)             & $\approx$5.9 \\
Sandbox compute            & 40--55 \\
\quad era total            & $\approx$545--560 \\
\midrule
\multicolumn{2}{@{}l}{\emph{Label-run era}} \\
Fixer API                  & 340.03 \\
GPU pod                    & $\approx$4 \\
Sandbox compute            & $\approx$25 \\
\quad era total            & $\approx$370 \\
\midrule
\multicolumn{2}{@{}l}{\emph{Training era}} \\
Fine-tuning run            & $\approx$86 \\
Corpus note generation     & $\approx$9 \\
\bottomrule
\end{tabular}

\end{table}


